% GENERATED FILE. Edit canonical lessons, then run npm run generate:books.
% canonical-source-hash: fnv1a64:ed3f50bfaa0fe7d5
% canonical-lessons: ES-C267-no-verbo, ES-C267-no-nada, ES-C267-pregunta-si-no, ES-C267-verdad-no, ES-C267-negar-preguntar

\chapter{Saying No, and Asking}
\label{ch:es-negate-and-ask}
\hlchaptermodality{267}

\begin{chapteropening}
\textbf{By the end of this chapter:} \emph{I can say that something is not so, and ask a question that can be answered yes or no.}

Hold a four-line exchange that denies something and then asks the other person to agree.
\end{chapteropening}

\section[no + verb]{no + verb — putting \emph{no} in front of the verb — the whole of Spanish negation}

\label{lesson:ES-C267-no-verbo}



\begin{quote}
One rule, and it is shorter than the English one.
\end{quote}

\begin{grammarlens}[title={Grammar lens: putting \emph{no} in front of the verb — the whole of Spanish negation}]
You already know \textbf{no}. You met it already as an answer: \emph{¿Hablas español?} — \emph{No.}

Here is the part nobody has told you yet, and it is the entire rule:

\begin{quote}
\textbf{Put \emph{no} immediately before the verb. That is negation. There is nothing else.}
\end{quote}

\emph{Hablo español} — I speak Spanish. \textbf{No} \emph{hablo español} — I do \textbf{not} speak Spanish.

English needs a helper — \emph{I do not speak} — and that \emph{do} is doing no work except carrying the negation. Spanish never borrowed the habit. The word goes in front of the verb and the sentence is finished.

That means every sentence you can already say, you can now also unsay.
\end{grammarlens}

\subsection*{Your turn}
Say these aloud:

\begin{itemize}
\raggedright
  \item three sentences you already know, then say each one again negated
  \item the same three as yes-or-no questions, letting your voice rise
\end{itemize}

\begin{tcolorbox}[breakable,colback=teal!4,colframe=teal!35!black,title={Before you move on}]
What goes in front of the verb to negate it? (\emph{\textbf{no}}.) Does Spanish need a helper word like English's \emph{do}? (\textbf{No — nothing else changes.})
\end{tcolorbox}
\section[no ... nada]{no ... nada — the double negative — where Spanish stacks what English forbids}

\label{lesson:ES-C267-no-nada}



\begin{quote}
One rule, and it is shorter than the English one.
\end{quote}

\begin{grammarlens}[title={Grammar lens: the double negative — where Spanish stacks what English forbids}]
English teachers spend years stamping out \emph{I don't see nothing}. Spanish requires it.

\emph{No veo} \textbf{nada} — literally \emph{I don't see nothing}, and it means \textbf{I see nothing}.

The rule: when a negative word comes AFTER the verb, \emph{no} still has to sit in front of it. Both halves stay.

\begin{quote}
\emph{No} veo \textbf{nada}.  ·  \emph{No} hablo con \textbf{nadie}.  ·  \emph{No} voy \textbf{nunca}.
\end{quote}

Move the negative word to the front and the \emph{no} disappears, because the sentence already has its negation:

\begin{quote}
\textbf{Nada} veo.  ·  \textbf{Nadie} habla.  ·  \textbf{Nunca} voy.
\end{quote}

Both are correct. The first is what people say.
\end{grammarlens}

\subsection*{Your turn}
Say these aloud:

\begin{itemize}
\raggedright
  \item three sentences you already know, then say each one again negated
  \item the same three as yes-or-no questions, letting your voice rise
\end{itemize}

\begin{tcolorbox}[breakable,colback=teal!4,colframe=teal!35!black,title={Before you move on}]
In \emph{no veo nada}, how many negative words are there? (\textbf{Two}, and both are required.) What happens to \emph{no} if \emph{nada} moves to the front? (\textbf{It goes away} — the sentence is already negative.)
\end{tcolorbox}
\section[¿Hablas español?]{¿Hablas español? — the yes-or-no question — and why Spanish barely has to do anything}

\label{lesson:ES-C267-pregunta-si-no}



\begin{quote}
One rule, and it is shorter than the English one.
\end{quote}

\begin{grammarlens}[title={Grammar lens: the yes-or-no question — and why Spanish barely has to do anything}]
English rebuilds the sentence to ask a yes-or-no question. \emph{You speak Spanish} becomes \emph{Do you speak Spanish?} — a word appears from nowhere and the order changes.

Spanish does neither. \textbf{The words stay exactly where they are.} Only your voice goes up, and the writing marks it:

\begin{quote}
\emph{Hablas español.}  $\to$  \textbf{¿Hablas español?}
\end{quote}

Same three words, same order. The \textbf{¿} at the front is the instruction: \emph{this is a question, start lifting your voice now.} That is why Spanish opens it — English readers do not find out until the final \emph{?}, by which time the sentence has been read flat.

You may also put the subject after the verb — \emph{¿Habla usted español?} — and with \emph{usted} that is the more usual shape. But it is a preference, not a rule. The rising voice is the question.
\end{grammarlens}

\subsection*{Your turn}
Say these aloud:

\begin{itemize}
\raggedright
  \item three sentences you already know, then say each one again negated
  \item the same three as yes-or-no questions, letting your voice rise
\end{itemize}

\begin{tcolorbox}[breakable,colback=teal!4,colframe=teal!35!black,title={Before you move on}]
What has to change to turn a Spanish statement into a yes-or-no question? (\textbf{Nothing but the voice} — and the ¿ ? marks.) Why does Spanish open with ¿? (\textbf{So the reader lifts their voice from the first word}, not the last.)
\end{tcolorbox}
\section[¿verdad?]{¿verdad? — the tag question — asking someone to agree}

\label{lesson:ES-C267-verdad-no}



\begin{quote}
One rule, and it is shorter than the English one.
\end{quote}

\begin{grammarlens}[title={Grammar lens: the tag question — asking someone to agree}]
There is a second kind of yes-or-no question: the one where you already think you know the answer and want the other person to agree.

English has a whole machine for this — \emph{isn't it, doesn't he, haven't you, aren't they} — and it changes with every verb. Spanish has two words.

\begin{quote}
\emph{Hablas español,} \textbf{¿verdad?} — \emph{You speak Spanish, right?} \emph{No hablas español,} \textbf{¿no?} — \emph{You don't speak Spanish, do you?}
\end{quote}

\textbf{¿verdad?} is \emph{truth?} and works everywhere. \textbf{¿no?} is quicker and slightly more casual. Neither one changes for the verb, the person, or the tense.

One tidy convention: after a negative sentence people prefer \textbf{¿verdad?}, since \emph{¿no?} after \emph{no} sounds like a stutter.
\end{grammarlens}

\subsection*{Your turn}
Say these aloud:

\begin{itemize}
\raggedright
  \item three sentences you already know, then say each one again negated
  \item the same three as yes-or-no questions, letting your voice rise
\end{itemize}

\begin{tcolorbox}[breakable,colback=teal!4,colframe=teal!35!black,title={Before you move on}]
What does Spanish use instead of English's \emph{isn't it / doesn't he / haven't you}? (\emph{\textbf{¿verdad?}} and \emph{\textbf{¿no?}}.) Do they change with the verb? (\textbf{Never.})
\end{tcolorbox}
\section[Practice]{negar y preguntar — saying something is not so, and asking whether it is}

\label{lesson:ES-C267-negar-preguntar}



\begin{quote}
Everything in this chapter, in one place.
\end{quote}

\begin{grammarlens}[title={Grammar lens: saying something is not so, and asking whether it is}]
This chapter gave you two things that go together more often than not.

\noindent
\begin{tabularx}{\linewidth}{@{}>{\raggedright\arraybackslash}X>{\raggedright\arraybackslash}X@{}}
\toprule
\textbf{you can now} & \textbf{how} \\
\midrule
say something is not so & \emph{no} in front of the verb \\
say it with a negative word after & \emph{no} stays, both halves \\
ask a yes-or-no question & nothing changes but your voice \\
ask someone to agree & \emph{¿verdad?} or \emph{¿no?} \\
\bottomrule
\end{tabularx}

Put them together and you have the shape most real conversations run on:

\begin{quote}
— \emph{¿Hablas español?} — \emph{No, no hablo español. Hablo inglés.} — \emph{Pero entiendes un poco,} \textbf{¿verdad?} — \emph{Sí, un poco.}
\end{quote}

Every word there is one you already had. What is new is that you can now \textbf{deny} and \textbf{ask}, which between them cover most of what a beginner needs to survive a conversation they did not start.
\end{grammarlens}

\subsection*{Your turn}
Say these aloud:

\begin{itemize}
\raggedright
  \item three sentences you already know, then say each one again negated
  \item the same three as yes-or-no questions, letting your voice rise
\end{itemize}

\begin{tcolorbox}[breakable,colback=teal!4,colframe=teal!35!black,title={Before you move on}]
How do you negate a verb? (\emph{\textbf{no}} in front of it.) How do you ask a yes-or-no question? (\textbf{Voice up, ¿ ?} — nothing else moves.) How do you ask someone to agree? (\emph{\textbf{¿verdad?}} or \emph{\textbf{¿no?}}.)
\end{tcolorbox}
